% Part: computability
% Chapter: computability-theory
% Section: russells-paradox

\documentclass[../../../include/open-logic-section]{subfiles}

\begin{document}

\olfileid{cmp}{thy}{rus}
\olsection{د رسل له تناقض سره پرتله}

ګټوره ده چې د دې برخې استدلالونه د رسل له تناقض سره پرتله او بېل
کړو:
\begin{enumerate}
\item د رسل تناقض: فرض کړئ $S = \Setabs{x}{x \notin x}$۔ بيا $S
  \in S$ هله او يوازې هله چې $S \notin S$ وي، چې تناقض دے۔

  \emph{پايله:} داسې سټ~$S$ نشته۔ د ``ټولو سټونو د سټ'' شتون منل د
  سټ تيورۍ له نورو بديهي اصولو سره ناسازګار دي۔
  \textbf{د ژباړې د نسخې سمون (OLCMP-023):} د سرچينې د دوه‌اړخيز
  شرط په دويم اړخ کښې ناڅاپه يو بل لوے توری راغلے و، که څه هم يوازې
  د سټ ټاکلے نوم تعريف شوے او د رسل شرط هم د هماغه سټ خپل غړيتوب
  ازمويي۔ دلته بل توری د ټاکلي سټ په نوم سم شوے دے۔

\item د رسل د تناقض يوه بدله بڼه: $F$ له ټولو تابعو له سټ څخه
  $\{ 0, 1 \}$ ته هغه ``تابع'' وګڼئ چې داسې تعريفېږي:
  \[
  F(f) =
  \begin{cases}
    1 & \text{که $f$ د~$f$ د تعريف ساحې غړے وي، او $f(f) = 0$ وي} \\
    0 & \text{که داسې نۀ وي}
  \end{cases}
  \]
  ورته استدلال ښيي چې $F(F) = 0$ هله او يوازې هله چې $F(F) = 1$
  وي، چې تناقض دے۔

  \emph{پايله:} $F$ تابع نۀ ده۔ د ``ټولو تابعو سټ'' دومره لوے دے
  چې د يوې تابع تعريف ساحه نۀ شي کېدے۔

\item قطري استدلال: $f_0$، $f_1$، \dots د جزوي محاسبه کېدونکو تابعو
  شمېرنه وي، او $G \colon \Nat \to \{ 0, 1 \}$ داسې تعريف کړئ:
  \[
  G(x) =
  \begin{cases}
    1 & \text{که $f_x(x)\downarrow = 0$ وي} \\
    0 & \text{که داسې نۀ وي}
  \end{cases}
  \]
  که $G$ محاسبه کېدونکې وي، نو د کوم~$k$ دپاره تابع~$f_k$ ده۔ خو
  بيا $G(k) = 1$ هله او يوازې هله چې $G(k) = 0$ وي، چې تناقض دے۔

  \emph{پايله:} $G$ محاسبه کېدونکې نۀ ده۔ پام وکړئ چې د سټ تيورۍ د
  بديهي اصولو له مخې~$G$ بيا هم تابع ده؛ دلته تناقض نشته، يوازې يوه
  څرګندونه ده۔
\end{enumerate}

د جزوي تابعو، محاسبه کېدونکو تابعو، جزوي محاسبه کېدونکو تابعو او
داسې نورو يادونه ګډوډوونکې کېدے شي۔ له~$\Nat$ څخه~$\Nat$ ته د ټولو
جزوي تابعو سټ د څيزونو يوه لويه ټولګه ده۔ ځينې ئې هرځاے تعريف شوې،
ځينې ئې محاسبه کېدونکې، ځينې دواړه هرځاے تعريف شوې او محاسبه
کېدونکې، او ځينې يوه هم نۀ دي۔ په ياد ولرئ چې کله بې له قيده
``تابع'' وايو، موخه مو هرځاے تعريف شوې تابع وي۔ نو دا ډولونه لرو:
\begin{enumerate}
\item محاسبه کېدونکې تابعې
\item هغه جزوي محاسبه کېدونکې تابعې چې هرځاے تعريف شوې نۀ دي
\item هغه تابعې چې محاسبه کېدونکې نۀ دي
\item هغه جزوي تابعې چې نۀ هرځاے تعريف شوې او نۀ محاسبه کېدونکې دي
\end{enumerate}
د دې د بېلولو دپاره ښايي داسې لوے مربع رسمول ګټور وي چې له~$\Nat$
څخه~$\Nat$ ته ټولې جزوي تابعې وښيي؛ بيا پکښې دوه يو پر بل ورغلي
سيمې په نښه کړئ، چې يوه د هرځاے تعريف شويو تابعو او بله د جزوي
محاسبه کېدونکو تابعو وي۔ ښه تمرين دے چې وګورئ د نقشې په هره جوړه شوې
سيمه کښې يو څيز بيانولے شئ که نه۔

\end{document}
